% !PRINTED 91
\weberpage{91}
%<<CONT
bleibt. Man erhält aber aus (12) einen Ausdruck, der auch für
ein reelles $a$ noch konvergiert, wenn man die Summe
\weq{13}{\sum_{1,\infty}^{\nu} e^{2\pi i\nu a} = -\frac{e^{2\pi ia}}{e^{2\pi ia} - 1}}
hinzufügt. Dadurch geht er nämlich über in
\[
\sum_{1,\infty}^{\nu} \frac{e^{2\pi i\nu a}q^{2\nu}e^{2\pi iv}}{q^{2\nu}e^{2\pi iv} - 1},
\]
und diese Reihe bleibt konvergent, solange der imaginäre Teil $a''$
von $a$ zwischen $-\omega''$ und $+\omega''$ liegt.

Demnach ergibt sich aus (7), wenn wir das dem Werte $\nu = 0$
entsprechende Glied absondern und die negativen $\nu$ durch $-\nu$
ersetzen:
\[
\frac{\vartheta'_{11}\vartheta_{11}(v+a)}{\pi\vartheta_{11}(v)\vartheta_{11}(a)} = \frac{2i}{e^{2\pi iv} - 1} + \frac{2ie^{2\pi ia}}{e^{2\pi ia} - 1}
\]
\[
+ 2i\sum_{1,\infty}^{\nu} \frac{e^{-2\pi i\nu a}}{q^{-2\nu}e^{2\pi iv} - 1} + 2i\sum_{1,\infty}^{\nu} \frac{q^{2\nu}e^{2\pi i\nu a}e^{2\pi iv}}{q^{2\nu}e^{2\pi iv} - 1}.
\]
Es ist aber
\[
\frac{2i}{e^{2\pi iv} - 1} = \frac{e^{-\pi iv}}{\sin\pi v} = \mathrm{cotg}\,\pi v - i,
\]
\[
\frac{2ie^{2\pi ia}}{e^{2\pi ia} - 1} = \frac{e^{\pi ia}}{\sin\pi a} = \mathrm{cotg}\,\pi a + i,
\]
und demnach läßt sich diese Entwickelung auch so darstellen:
\weq{14}{\frac{\vartheta'_{11}\vartheta_{11}(v+a)}{\pi\vartheta_{11}(v)\vartheta_{11}(a)} = \mathrm{cotg}\,\pi v + \mathrm{cotg}\,\pi a
-2i\sum \left(\frac{q^m e^{m\pi ia}}{e^{-2\pi iv} - q^m} - \frac{q^m e^{-m\pi ia}}{e^{2\pi iv} - q^m}\right),}
worin $m$ die \gesperrt{Reihe der geraden Zahlen}
\[
m = 2, 4, 6, 8, \dots
\]
durchläuft. Der Gültigkeitsbereich dieser Entwickelung ist
\[
-\omega'' < a'' < \omega''.
\]
In der Tabelle II sind diese 16 Entwickelungen zusammengestellt.

Aus diesen Formeln sind drei andere ableitbar, indem man
$v$ oder $a$ oder beide um $\frac{1}{2}$ vermehrt. Drei andere aber, die der
Vermehrung von $v$ und $a$ um $\frac{1}{2}\omega$ entsprechen, müssen direkt abgeleitet
 werden. Sie zeigen reelle Form, wenn $q$, $a$, $v$ reell sind.
